Las instrucciones de uso del software se detallan a continuación. Este manual se dirige al usuario final del sistema \textbf{UniHub}.

\section{Introducción}
El Portal Web de \textbf{UniHub} (Fase 3) permite consultar el catálogo oficial de universidades de España, explorar titulaciones de Grado, Máster y Doctorado y, cuando exista una fuente curricular verificada, revisar su desglose de asignaturas, créditos ECTS y enlace al documento oficial.

\section{Instalación}
Para acceder al sistema no se requiere ninguna instalación por parte del usuario final. Únicamente se necesita un navegador web moderno (Google Chrome, Mozilla Firefox, Microsoft Edge o Safari) y conexión a la red donde esté desplegado el servicio (\url{http://localhost} o \url{http://localhost:5173}).

\section{Uso del sistema}
\begin{enumerate}
    \item \textbf{Exploración de Universidades:} Acceda a la pestaña "Universidades" para filtrar centros públicos (listados prioritariamente en primer lugar) o privados por nombre o comunidad autónoma.
    \item \textbf{Buscador de Titulaciones y BOE:} En la pestaña "Titulaciones", busque por cualquier titulación o filtre por Grado, Máster o Doctorado. Al abrir una tarjeta se muestra la estructura curricular disponible: para Grados y Másteres se desglosan las asignaturas organizadas por curso/cuatrimestre y menciones curriculares; para Doctorados (RD 99/2011) se visualiza la insignia oficial verificada, la Escuela de Doctorado tutelar, la cuadrícula de líneas de investigación científica acreditadas, actividades formativas transversales y el régimen de tutela académica anual, junto con el enlace directo a la fuente oficial. Algunos documentos oficiales sólo publican un resumen de créditos; en esos casos el portal no inventa asignaturas ausentes.
    \item \textbf{Paginación Configurable:} En la parte inferior de las grillas de universidades y titulaciones, seleccione el número de elementos a mostrar por página (5, 10, 20, 50 o 100) y navegue entre páginas con los botones de control.
    \item \textbf{Búsqueda por Cercanía (Geolocalización):} En la pestaña "Por Cercanía", pulse el botón "Usar Mi Ubicación Actual (GPS)" para permitir que el navegador calcule la distancia exacta en km a cada universidad y visualícela directamente dentro de la tarjeta de cada centro.
    \item \textbf{Sección "Sobre Nosotros":} Acceda a la pestaña "Sobre Nosotros" para consultar la información institucional del proyecto (Trabajo Fin de Grado de Alejandro Ramos Rodríguez en la Universidad de Cádiz).
    \item \textbf{Panel del Administrador (Acceso Aislado):} Escriba la ruta manual \url{http://localhost/admin} en la barra de direcciones de su navegador. Introduzca la API Key definida en \texttt{ADMIN\_API\_KEY} para acceder al panel de administración. No se distribuyen credenciales por defecto.
\end{enumerate}
